%==============================================================
%                    CAPITULO 1
%==============================================================
\chapter{Fundamentos teóricos del problema de planificación de rutas y su solución mediante diferentes métodos}\label{cap:fundamentos}

En este capítulo se describen los aspectos teóricos de la investigación que se encuentran relacionados con los MDVRP. Entre los aspectos que se abordan están las características y variantes del problema, así como los diferentes métodos de solución encontrados en la literatura. Además, se hace énfasis en los algoritmos de agrupamiento como alternativa para la fase de asignación de clientes a depósitos. Por último, se describe la Biblioteca de Heurísticas de Asignación para Problemas de Planificación de Rutas de Vehículos con Múltiples Depósitos y sus algoritmos existentes.

\section{Problemas de Planificación de Rutas de Vehículos}

El concepto de Problema de Planificación de Rutas de Vehículos (\textit{Vehicle Routing Problem}, VRP) es una generalización del Problema del Viajero (\textit{Traveling Salesman Problem}, TSP). En concreto, el VRP es un problema de optimización combinatoria cuyo objetivo es atender las demandas de clientes que se encuentran dispersos geográficamente sujetos a un conjunto de restricciones. La diferencia radica en que TSP dispone de un mapa con múltiples puntos que un individuo tiene que visitar para volver a su punto de origen realizando la ruta más corta posible. En cambio, el VRP no se limita a solo una persona, sino que pueden existir más de un individuo, en este caso vehículos, que normalmente deben entregar objetos en diversos puntos de la ciudad a clientes. De este modo, el problema se vuelve más complejo que el TSP \parencite{kumarSurveyVehicleRouting2012}.

Las aplicaciones informáticas para el dise\~no de rutas se derivan de modelos VRP que atienden diferentes necesidades e incluyen diversas restricciones. Esto ha dado como resultado el surgimiento de diferentes variantes del clásico VRP, de las cuales se pueden mencionar las siguientes:

\begin{itemize}
  \item VRP con recolección y entrega (VRPPD): estudia el caso donde una empresa debe recolectar y entregar bienes en cantidades específicas para cada cliente visitado \parencite{berbegliaDynamicPickupDelivery2007}.
  \item VRP con flota heterogénea: los vehículos tienen diferentes capacidades de carga  \parencite{baldacciRoutingHeterogeneousFleet2008}.
  \item \textit{Open VRP} (OVRP): se refiere al caso de cuando el vehículo no regresa al depósito de origen, es decir, que se dirige a otro depósito o sitio diferente para cargar productos, en un contexto de flota rentada \parencite{sariklisHeuristicMethodOpen2000}.
  \item VRP con ventanas de tiempo (VRPTW): plantea que cada cliente tiene que ser atendido de manera obligada dentro de un cierto horario específico \parencite{cordeauVRPTimeWindows2002}.
  \item VRP con depósitos múltiples (MDVRP): implica que la empresa posee diversos depósitos desde los cuales puede abastecer a los clientes \parencite{montoya-torresLiteratureReviewVehicle2015}.
  \item \textit{Stochastic VRP} (SVRP): Propone que las variables involucradas son desconocidas o aleatorias, e incluso que se pueden encontrar en un rango de probabilidad, tales como el número de clientes, sus demandas, tiempo de servicio, tiempo de viaje, entre otras \parencite{gendreauStochasticVehicleRouting1996}.
  \item \textit{Periodic VRP} (PVRP): analiza un problema donde los pedidos pueden ser llevados sólo en ciertos días  \parencite{francisPeriodVehicleRouting2008}.
  \item Capacitated VRP: estudia el caso donde los vehículos tienen una capacidad máxima de carga y se deben dise\~nar rutas que satisfagan la demanda de los clientes sin exceder dicha capacidad \parencite{ibrahimCAPACITATEDVEHICLEROUTING2019}.
\end{itemize}

\subsection{Problema de Planificación de Rutas de Vehículos con Múltiples Depósitos}

El MDVRP es un problema de optimización combinatoria, categorizado por su complejidad como NP-Duro  \parencite{rebillasloredoMultidepotVRPVehicle2018}. Por lo que es difícil encontrar su solución óptima usando métodos exactos en un tiempo polinomial. Esta variante cuenta con varios depósitos y cada uno tiene su propia flota de vehículos  \parencite{colinDisenoAlgoritmoMetaheuristico2024}.

El problema puede ser representado mediante modelos de optimización y se basa en  \parencite{montoya-torresLiteratureReviewVehicle2015}, \parencite{Cordeau2012}, \parencite{Vidal2012}, \parencite{Subramanian2013}, \parencite{Escobar2014}, \parencite{cordeauTabuSearchHeuristic1997}. Sea \(G = (V, A)\) un grafo completo, donde \(V\) es el conjunto de nodos y \(A\) es el conjunto de arcos. Los nodos se dividen en dos subconjuntos: los clientes a ser atendidos, \(V_{C} = \{1, \ldots, N\}\), y los múltiples depósitos \(V_{D} = \{N + 1, \ldots, N + M\}\), con \(V_{C} \cup V_{D} = V\) y \(V_{C} \cap V_{D} = \emptyset\). Hay un costo no negativo \(c_{ij}\) asociado con cada arco \((i, j) \in A\). La demanda de cada cliente es \(d_{i}\) donde no hay demanda en los nodos de depósito. También hay una flota de \(K\) vehículos, cada uno con capacidad \(Q\). El tiempo de servicio en cada cliente \(i\) es \(t_{i}\) mientras que el tiempo máximo de duración de la ruta se establece en \(T\). Es posible que se necesite un factor de conversión \(w_{ij}\) para transformar el costo \(c_{ij}\) en unidades de tiempo. Sin embargo, el costo es el mismo que las unidades de tiempo y distancia, por lo que \(w_{ij} = 1\)  \parencite{oliveiraCooperativeCoevolutionaryModels2015}.

A continuación se presenta la formulación matemática \parencite{oliveiraCooperativeCoevolutionaryModels2015}:

Las variables binarias \(x_{ijk}\) son iguales a 1 cuando el vehículo \(k\) visita el nodo \(j\) inmediatamente después del nodo \(i\). Las variables auxiliares \(y_{i}\) también se utilizan en las restricciones de eliminación de subtours.

\noindent\textbf{Función Objetivo:} Minimizar el costo total.
\begin{equation}
\mathrm{Minimize}\sum_{i = 1}^{N + M}\sum_{j = 1}^{N + M}\sum_{k = 1}^{K}c_{ij}x_{ijk},
\end{equation}

\noindent\textbf{Restricciones:}

Garantiza que cada cliente sea visitado exactamente una vez (flujo entrante).
\begin{equation}
\sum_{i = 1}^{N + M}\sum_{k = 1}^{K}x_{ijk} = 1 \quad (j = 1,\dots ,N);
\end{equation}

Garantiza que cada cliente sea visitado exactamente una vez (flujo saliente).
\begin{equation}
\sum_{j = 1}^{N + M}\sum_{k = 1}^{K}x_{ijk} = 1 \quad (i = 1,\dots ,N);
\end{equation}

Conservación de flujo. Lo que entra a un nodo es igual a lo que sale de él para cada vehículo.
\begin{equation}
\sum_{i = 1}^{N + M}x_{ihk} - \sum_{j = 1}^{N + M}x_{hjk} = 0 \quad (k = 1,\dots ,K;h = 1,\dots ,N + M);
\end{equation}

Restricción de capacidad. La demanda total de la ruta no debe exceder la capacidad del vehículo \(Q\).
\begin{equation}
\sum_{i = 1}^{N + M}\sum_{j = 1}^{N + M}d_{i}x_{ijk}\leq Q \quad (k = 1,\dots ,K);
\end{equation}

Restricción de duración de la ruta. El costo total más el tiempo de servicio no debe exceder el tiempo máximo \(T\).
\begin{equation}
\sum_{i = 1}^{N + M}\sum_{j = 1}^{N + M}(c_{ij}w_{ij} + t_i)x_{ijk}\leq T \quad (k = 1,\dots,K);
\end{equation}

Verifica la disponibilidad de vehículos (salida del depósito).
\begin{equation}
\sum_{i = N + 1}^{N + M}\sum_{j = 1}^{N}x_{ijk}\leq 1 \quad (k = 1,\dots,K);
\end{equation}

Verifica la disponibilidad de vehículos (entrada al depósito).
\begin{equation}
\sum_{j = N + 1}^{N + M}\sum_{i = 1}^{N}x_{ijk}\leq 1 \quad (k = 1,\dots,K);
\end{equation}

Restricciones de eliminación de subtours.
\begin{equation}
y_{i} - y_{j} + (M + N)x_{ijk}\leq N + M - 1; \quad \mathrm{para} \ 1\leq i\neq j\leq N \ \mathrm{y} \ 1\leq k\leq K;
\end{equation}

Define la naturaleza binaria de las variables de decisión de ruta.
\begin{equation}
x_{ijk}\in \{0,1\} \quad \forall i,j,k;
\end{equation}

Define la naturaleza binaria de las variables auxiliares de eliminación de subtours.
\begin{equation}
y_{i}\in \{0,1\} \quad \forall i;
\end{equation}

En la Figura~\ref{fig:mdvrp} se ilustra la estructura del problema:

\begin{itemize}
  \item Se observan tres depósitos (representados por cuadrados), cada uno con una flota de vehículos.
  \item Los clientes (círculos rojos) son asignados a un depósito específico.
  \item Las líneas verdes muestran las rutas de los vehículos, que parten y regresan a su depósito de origen.
\end{itemize}

\figura{fig01_mdvrp.png}{0.75\textwidth}{Representación gráfica de MDVRP.}{fig:mdvrp}

La literatura documenta múltiples aplicaciones del MDVRP a problemas logísticos reales. Un estudio aplicado a una red de distribución de productos de limpieza y cuidado personal en Turquía utiliza el MDVRP para optimizar las rutas desde tres depósitos a 502 clientes. Mediante un enfoque de agrupar primero y planificar después se logran mejoras significativas en la eficiencia de la distribución, utilizando datos reales de demanda y coordenadas geográficas \parencite{hamidOptimizationEssentialsTheory2024}.

Otro ejemplo se encuentra en el Reino Unido, donde se desarrolla un modelo de programación entera mixta (MIP) \parencite{junger50YearsInteger2009} para resolver un MDVRP que incluye complejidades del mundo real como una flota heterogénea, costos de conductores y ventanas de tiempo. La implementación de este modelo permite reducir sustancialmente los costos de transporte sin necesidad de abrir nuevas instalaciones \parencite{mogaleModellingOptimisingMultidepot2025}.

También se encuentra un caso de estudio en Estambul que trata sobre la distribución de una empresa de bebidas utilizando datos históricos de GPS para construir una matriz de tiempo-distancia más precisa. Se evalúan múltiples enfoques de optimización, incluyendo algoritmos evolutivos y de aprendizaje profundo, para resolver el MDVRP. Los resultados muestran que estos métodos son efectivos para manejar la complejidad y la naturaleza dinámica de las entregas en entornos urbanos de alta densidad  \parencite{erdemAIInfusedEvolutionaryOptimization2025}.

\section{Algoritmos heurísticos}

Debido a la complejidad que presentan los VRP, existen dos tipos de algoritmos que pueden resolver este tipo de problemas, los métodos exactos y aproximados. Cada uno de ellos posee diferentes características que pueden considerarse ventajas y/o desventajas que deben ser tomadas en cuenta al ser utilizados \parencite{colinDisenoAlgoritmoMetaheuristico2024}. En la Figura~\ref{fig:tecnicas} se puede apreciar una taxonomía de estos métodos:

\figura{fig02_tecnicas.png}{0.85\textwidth}{Técnicas utilizadas para la solución de MDVRP \parencite{talbiMetaheuristicsDesignImplementation2009}.}{fig:tecnicas}

Entre los métodos exactos se pueden distinguir tres grupos: búsqueda directa de árbol, programación dinámica y programación lineal, entera y mixta. Todos estos métodos son eficientes en problemas de hasta 50 clientes, debido a restricciones de tiempo computacional. Se han llevado a cabo algunas investigaciones introduciendo métodos exactos para MDVRP  \parencite{colinDisenoAlgoritmoMetaheuristico2024}. Algunos ejemplos de los más utilizados para el MDVRP de este tipo de métodos son:

\begin{itemize}
  \item Ramificación y Corte: los algoritmos \textit{Branch \& Cut} se basan en el cálculo simultáneo de cotas inferiores y superiores de la solución óptima. Las cotas superiores se obtienen por medio de la resolución de las relajaciones lineales asociadas a cada nodo del árbol de enumeración, eventualmente ajustados por medio de planos de corte, mientras que las cotas inferiores son halladas por medio de procedimientos heurísticos y cuando la relajación lineal asociada con algún subproblema tiene óptimo entero  \parencite{hernandezAPLICACIONMODELOVRP2015}.
  \item Programación Lineal: la programación lineal es un método a través del cual se optimiza una función objetivo, bien sea maximizando o minimizando dicha función, en la cual las variables están elevadas a la potencia 1. Este campo de la programación lineal es un área de la programación matemática, dedicada a maximizar o minimizar una función lineal que recibe el nombre de ``función objetivo'', de manera que las variables de tal función estén sujetas a una serie de restricciones expresadas a través de un sistema de ecuaciones o inecuaciones lineales \parencite{colinDisenoAlgoritmoMetaheuristico2024}.
\end{itemize}

La investigación se centra en los métodos heurísticos, los cuales se han caracterizado por ser procedimientos que garantizan buenas soluciones mediante una exploración limitada del espacio de búsqueda. Los métodos heurísticos se pueden clasificar en métodos constructivos y heurísticas de mejora \parencite{colinDisenoAlgoritmoMetaheuristico2024}. Algunas heurísticas utilizadas en la literatura para solucionar el MDVRP son:

\begin{itemize}
  \item Heurística ad hoc: en 1993 fue presentada una nueva heurística ad hoc en el aporte de \parencite{chaoNewHeuristicMultidepot1993}, su método de solución asigna los clientes a los depósitos de manera fácil y rápida. En el procedimiento de mejora, permite que la distancia total recorrida por la flota aumente con la esperanza de encontrar finalmente rutas muy cortas \parencite{colinDisenoAlgoritmoMetaheuristico2024}.
  \item Algoritmo de Ahorros: extensión del método clásico para VRP al contexto de múltiples depósitos. Asigna clientes a depósitos y luego construye rutas fusionando recorridos según el ahorro de distancia \parencite{clarkeSchedulingVehiclesCentral1964}. En  \parencite{tillmanMultipleTerminalDelivery1969} se presentó la primera adaptación para MDVRP.
  \item Algoritmo de Barrido: método de dos fases: primero agrupa clientes por ángulo polar respecto al depósito, luego construye rutas para cada grupo. Bueno para problemas de tama\~no medio-grande \parencite{gillettHeuristicAlgorithmVehicledispatch1974}.
  \item \textit{Merge-Head and Drop-Tail}: una técnica planteada para resolver el MDVRP es la denominada estrategia \textit{Merge-Head and Drop-Tail} (MHDT), combinada con optimización de combinación de sectores (SCO), aplicada por  \parencite{shiHeuristicSolutionMethod2020}; aquí se plantea un comportamiento inverso a la entrega, es decir debe generar rutas de recolección, dada esta única diferencia, el resto de las características y restricciones son idénticas, su propuesta funciona de igual manera para el MDVRP. En este método, primero asignan puntos de recolección de acuerdo con la distancia más cercana a los puntos de recolección, luego cada planta resuelve el VRP respectivamente, asignando clientes a los vehículos y planificando el orden en que los vehículos visitan a los clientes. En el último paso, proponen el algoritmo de optimización de combinación de sectores (SCO) para generar múltiples soluciones iniciales, y luego estas soluciones iniciales se mejoran utilizando la estrategia de combinación de cabeza y caída (MHDT).
\end{itemize}

Las metaheurísticas son estrategias de alto nivel que guían y modifican heurísticas subyacentes para explorar el espacio de búsqueda de forma eficiente. Estos algoritmos se basan en la exploración iterativa del espacio de soluciones, utilizando estrategias para escapar de los óptimos locales. Estas se clasifican en basadas en población o basadas en un punto  \parencite{blumMetaheuristicsCombinatorialOptimization2003}.

La Búsqueda Tabú se mantiene como uno de los mejores enfoques para resolver el MDVRP. Esta es una metaheurística que explora el espacio de soluciones permitiendo movimientos que empeoran la solución actual para escapar de óptimos locales. Para evitar ciclos, mantiene una lista de movimientos prohibidos o tabú durante un número determinado de iteraciones  \parencite{hernandezAPLICACIONMODELOVRP2015}.

También otra de las metaheurísticas de punto único que se destaca es la Búsqueda Local Iterada. Este algoritmo parte de una solución inicial, aplica una perturbación para diversificar la búsqueda, luego una búsqueda local para alcanzar un nuevo óptimo y finalmente decide si acepta la nueva solución según algún criterio \parencite{colinDisenoAlgoritmoMetaheuristico2024}.

Entre los poblacionales se destacan bastante los genéticos por la flexibilidad ante ciertas restricciones. Estos algoritmos mantienen una población de soluciones que evoluciona mediante operadores de selección, cruce y mutación. Para el MDVRP, las soluciones se representan como secuencias de clientes. El cruce combina sub-rutas de dos padres para generar un descendiente, reparando la factibilidad si es necesario \parencite{hernandezAPLICACIONMODELOVRP2015}.

Otro de los algoritmos poblacionales que se destacan fue el Algoritmo Colonia de Hormigas. Esta simula el comportamiento de hormigas reales que depositan feromona en el camino. Las hormigas artificiales construyen soluciones probabilísticamente guiadas por la feromona (que refleja la calidad de soluciones previas) y por información heurística local. Para MDVRP, se optimizan rutas por depósito y se a\~naden operadores de cruce para mejorar el rendimiento \parencite{colinDisenoAlgoritmoMetaheuristico2024}.

\section{Algoritmos de agrupamiento para la fase de asignación de clientes a depósitos}

La minería de datos es un campo interdisciplinario de la computación que se centra en el descubrimiento de patrones significativos en grandes volúmenes de datos mediante el uso de métodos avanzados de inteligencia artificial, aprendizaje automático, estadística y bases de datos. Su objetivo principal es extraer conocimiento valioso de los datos y transformarlo en una estructura comprensible, ya sea para su análisis o para su almacenamiento en bases de datos con miras a usos posteriores \parencite{oralloIntroduccionMineriaDatos2004}.

La minería de datos se basa en una variedad de técnicas y métodos que permiten extraer información útil de grandes volúmenes de datos. Algunas de las técnicas más comunes incluyen \parencite{kesavarajStudyClassificationTechniques2013}:

\begin{itemize}
  \item Clasificación: es una técnica utilizada para asignar elementos de un conjunto de datos a categorías predefinidas. Esto se logra mediante algoritmos de aprendizaje supervisado, como los árboles de decisión, las máquinas de vectores de soporte (SVM) y las redes neuronales.
  \item Agrupamiento: es una técnica no supervisada que organiza los datos en clústeres basados en la similitud entre los datos. Esta no requiere categorías predefinidas.
  \item Regresión: es una técnica utilizada para predecir un valor numérico basado en la relación entre variables.
  \item Reglas de asociación: identifican relaciones entre variables en un conjunto de datos.
\end{itemize}

La fase de asignación de clientes a depósitos en MDVRP puede considerarse un problema de agrupamiento, ya que el resultado es un conjunto de clientes agrupados por depósito. Esta interpretación sugiere un enfoque que agrupa a los clientes y luego programa los vehículos en cada grupo. En la literatura el MDVRP puede resolverse en dos etapas: primero, los clientes deben asignarse a los depósitos; luego, se deben construir rutas que vinculen a los clientes asignados al mismo depósito. Un enfoque razonable sería dividir el problema en tantos subproblemas como depósitos existan y resolver cada subproblema por separado, lo cual transforma el MDVRP en múltiples VRP independientes \parencite{sekerSolvingMultiDepot2009}. Por tanto, resulta interesante adaptar los algoritmos de agrupamiento para la fase de asignación en el MDVRP.

Los algoritmos de agrupamiento se basan en la cuantificación de qué tan similares o no sean los objetos del conjunto de datos. Las métricas de distancia desempe\~nan un papel central en la agrupación de estos objetos en grupos. Entre las medidas comunes de similitud y distancia se incluyen varios tipos de distancia como \textit{Euclidian}, \textit{Manhattan}, \textit{Haversine} y \textit{Chebyshev}. La elección de la medida de proximidad influye en los resultados de la agrupación, ya que los distintos conjuntos de datos requieren técnicas de separación específicas y métricas de similitud adaptadas a sus características. Estos algoritmos de agrupamiento según  \parencite{koutroumbasPatternRecognition2008}, se pueden clasificar como:

\begin{itemize}
  \item Agrupamiento particional: en estos algoritmos cada partición representa un grupo y cada objeto debe pertenecer exactamente a un grupo. Si se aplica la partición difusa, permite que los objetos pertenezcan a múltiples grupos con probabilidades. Aquí se encuentran algoritmos como k-means, k-medoids, k-modes y c-means difusos. Estos métodos son sensibles a la inicialización y a los valores atípicos, y se obtienen resultados confiables cuando los grupos son esféricos, de tama\~no y distribución uniformes. En la Figura~\ref{fig:particional} se puede apreciar gráficamente el funcionamiento de este tipo de agrupamientos.
  \item Agrupamiento jerárquico: se dividen en enfoques basados en aglomeración (de abajo hacia arriba) y divisivos (de arriba hacia abajo), que producen una jerarquía de grupos representados como dendrogramas. Los métodos basados en aglomeración comienzan con cada punto de datos como un grupo separado y fusionan los grupos más cercanos de forma iterativa, mientras que los métodos divisivos comienzan con todos los puntos de datos en un grupo y los dividen de forma recursiva. Los criterios de vinculación incluyen la vinculación simple, completa, promedio y el método de Ward, cada uno de los cuales define la proximidad del grupo de manera diferente. Los ejemplos de algoritmos de agrupamiento jerárquico incluyen BIRCH, CURE y ROCK. En la Figura~\ref{fig:jerarquico} se puede apreciar gráficamente el funcionamiento de este tipo de agrupamientos.
  \item Agrupamiento basado en densidad: los algoritmos de agrupamiento basados en densidad identifican los agrupamientos como regiones de alta densidad separadas por áreas de baja densidad, sin requerir el número de agrupamientos como entrada. Dentro de esta clasificación se encuentra el agrupamiento espacial basado en densidad de aplicaciones con ruido (DBSCAN), el agrupamiento de puntos para identificar la estructura del agrupamiento (OPTICS), el agrupamiento de picos de densidad (DPC) y el agrupamiento jerárquico basado en densidad (HDBSCAN), este último destaca por construir una jerarquía de clústeres que le permite manejar densidades variables sin requerir un parámetro global de radio. Estos métodos pueden presentar dificultades con densidades variables. En la Figura~\ref{fig:densidad} se puede apreciar gráficamente el funcionamiento de este tipo de agrupamientos.
\end{itemize}

\figura{fig03_particional.png}{1\textwidth}{Representación gráfica de algoritmos de agrupamiento particionales \parencite{gaoOverviewClusteringMethods2023}.}{fig:particional}

\figura{fig04_jerarquico.png}{1\textwidth}{Representación gráfica de algoritmos de agrupamiento jerárquicos \parencite{gaoOverviewClusteringMethods2023}.}{fig:jerarquico}

\figura{fig05_densidad.png}{1\textwidth}{Representación gráfica de algoritmos de agrupamiento basados en densidad \parencite{gaoOverviewClusteringMethods2023}.}{fig:densidad}

Los algoritmos de agrupamiento se aplican al MDVRP principalmente bajo el enfoque ``agrupar primero, planificar después'', donde primero se agrupan clientes y luego se construyen rutas para cada grupo. Un buen ejemplo de la aplicación de esto es una distribución de alimentos a supermercados en Reino Unido donde se utilizó K-means para agrupar 120 supermercados en cuatro clústeres correspondientes a los depósitos disponibles, minimizando la distancia intra-cluster. Luego se optimizaron rutas independientes para cada clúster. El ahorro en costos operativos fue del 15\% \parencite{salhiHeuristicAlgorithmMultidepot1998}.

\section{Biblioteca de Heurísticas de Asignación en la Planificación de Rutas de Vehículos con Múltiples Depósitos}

El área de optimización ha llevado a la creación de diversas bibliotecas en diferentes lenguajes de programación, cada una con su propia variedad de algoritmos y características. A continuación, se presenta una lista no exhaustiva de bibliotecas destacadas en resolver VRP y sus variantes, entre las cuales se incluye el MDVRP. De cada una se ofrece una breve descripción de sus características y usos  \parencite{ericramosaragonIncorporacionModuloMetricas2025}:

\begin{itemize}
  \item JSprit: es una biblioteca de código abierto escrita en Java. Es un componente ligero y flexible que resulta fácil de utilizar, fundamentado en una única metaheurística multipropósito que actualmente aborda una variedad amplia de problemas, como el Problema del Agente Viajero y sus variantes de VRP. Aunque no está específicamente dise\~nada para abordar la fase de asignación en problemas con múltiples depósitos, JSprit puede implementar estrategias de agrupamiento geográfico como un paso preliminar, utilizando técnicas como ordenamientos angulares o agrupaciones espaciales similares al enfoque del clásico algoritmo de barrido. Estas estrategias permiten organizar a los clientes en función de su proximidad antes de aplicar su metaheurística interna  \parencite{jspritcontributorsJSpritJavabasedLibrary2026}.
  \item \textit{Google Operation Research Tools} (OR-Tools): es un software de código abierto desarrollado en C++, con versiones estables disponibles en Java, C\# y Python. Originalmente dise\~nado para abordar diversos problemas de optimización como la programación lineal, la programación entera y la programación con restricciones, esta herramienta ha ido incorporando gradualmente capacidades para resolver variantes del VRP. Aunque brinda herramientas para construir modelos completos de VRP, su arquitectura flexible permite personalizar fases específicas, como la asignación, mediante programación y heurísticas personalizadas. Además, se pueden integrar algoritmos de agrupamiento para realizar una asignación previa de clientes a los depósitos. También es posible implementar restricciones a medida en el modelo, como límites de capacidad o balanceo de carga, para ajustar la asignación inicial antes de continuar con la planificación de las rutas de los vehículos. Esta herramienta resuelve varias variantes de VRP de manera nativa, pero no el MDVRP \parencite{perronORTools}.
  \item \textit{Vehicle Routing Open-Source Optimization Machine} (VROOM): es un componente de software para problemas de planificación de rutas de vehículos desarrollado en el lenguaje de programación C++. Esta herramienta utiliza \textit{Open Street Map} (OSM) y \textit{Open Source Routing Machine} (OSRM) como \textit{backend} para obtener rutas y devolver soluciones para las diferentes variantes del VRP. VROOM puede resolver distintas variantes del VRP que incluyen múltiples depósitos, capacidades, ventanas de tiempo, así como entrega y recogida \parencite{karkulaComparisonCapabilitiesRecent2019}.
  \item Scikit-learn: es una biblioteca de aprendizaje automático en Python que, aunque no está dise\~nada específicamente para resolver problemas de optimización como el MDVRP, proporciona una amplia gama de algoritmos de agrupamiento y análisis de datos que pueden ser utilizados para abordar la fase de asignación en este tipo de problemas. Algoritmos como K-means, DBSCAN y el agrupamiento jerárquico son útiles para preprocesar los datos y asignar clientes a depósitos de manera eficiente, basándose en proximidad y restricciones de capacidad. Su flexibilidad y compatibilidad con otras bibliotecas de Python la convierten en una herramienta versátil para integrar técnicas de agrupamiento en soluciones personalizadas para el MDVRP  \parencite{f.pedregosaScikitlearnMachineLearning2011}.
\end{itemize}

En el ámbito nacional, se han desarrollado diversas soluciones en el contexto académico orientadas a la resolución de problemas de planificación de rutas de vehículos, como parte del proyecto de Optimización y Metaheurísticas de la Facultad de Ingeniería Informática de la Universidad Tecnológica de La Habana ``José Antonio Echeverría'' (CUJAE). Algunas de estas herramientas se caracterizan por su enfoque modular y por estar dise\~nadas para abordar distintas variantes del VRP, con especial énfasis en la utilización de heurísticas y metaheurísticas aplicadas al MDVRP \parencite{ericramosaragonIncorporacionModuloMetricas2025}:

\begin{itemize}
  \item Bibliotecas de Heurísticas de Construcción para Problemas de Planificación de Rutas de Vehículos (BHCVRP): es un componente que implementa diez heurísticas de construcción para resolver siete variantes del VRP, incluyendo el MDVRP. Actualmente las soluciones obtenidas mediante estas heurísticas pueden ser utilizadas como solución final o como parte de una solución más compleja en un contexto donde se empleen metaheurísticas \parencite{anandadelacaridadmoralesmoralesNuevaVersionBiblioteca2025}.
  \item Capa de servicio para soluciones a Problemas de Planificación de Rutas de Vehículos (CaSVRP): es una capa de servicios web dirigidos a la resolución de diferentes variantes de VRP mediante algoritmos metaheurísticos. Entre las variantes del VRP que resuelve se halla MDVRP \parencite{g.yanezDesarrolloNuevosModelos2021}. Esta solución se compone de un módulo de optimización donde se utiliza la Biblioteca de Clases para la Integración de Algoritmos Metaheurísticos (BiCIAM)\parencite{hernandezNuevaVersionBiblioteca2014a}, un módulo para la generación de soluciones iniciales empleando BHCVRP y otro para la generación de nuevas soluciones a partir de otras soluciones haciendo uso de LMOVRP \parencite{loresOperadoresMutacionBasados2017}.
  \item Biblioteca de Heurísticas de Asignación para Problemas de Planificación de Rutas de Vehículos con Múltiples Depósitos (BHAVRP): es una biblioteca que resuelve la fase de asignación de clientes a depósitos en el MDVRP. Esta herramienta contiene 20 algoritmos de agrupamiento y heurísticas clásicas de asignación, y es utilizada por otras herramientas como BHCVRP, CaSVRP y TransO \parencite{ericramosaragonIncorporacionModuloMetricas2025}.
\end{itemize}

La Tabla~\ref{tab:bibliotecas} presenta una comparación entre las bibliotecas antes mencionadas que resuelven la fase de asignación del MDVRP. Es importante resaltar criterios como el hecho de ser de código abierto ya que eso sería de utilidad para seguir mejorando el componente.

\begin{table}[H]
\centering
\caption{Comparación entre las herramientas que resuelven el MDVRP.}
\label{tab:bibliotecas}
\small
\begin{tabular}{@{}lcccc@{}}
\toprule
\textbf{Biblioteca} & \textbf{\shortstack{Implementan\\heurísticas clásicas}} & \textbf{\shortstack{Implementan\\algoritmos de\\agrupamiento}} & \textbf{\shortstack{Lenguajes\\soportados}} & \textbf{\shortstack{Código\\abierto}} \\
\midrule
JSprit & No & No & Java & Sí \\
OR-Tools & Sí & No & \shortstack{C++, Python,\\Java, C\#} & Sí \\
VROOM & No & No & C++ & No \\
Scikit-learn & No & Sí & Python & Sí \\
CaSVRP & Sí & No & Java & Sí \\
BHCVRP & Sí & No & \shortstack{Java, Python} & Sí \\
BHAVRP & Sí & Sí & \shortstack{Java, Python} & Sí \\
\bottomrule
\end{tabular}
\end{table}

A partir del análisis comparativo de las bibliotecas podemos ver que solo BHAVRP y Scikit-learn ofrecen algoritmos de agrupamiento, lo que los posiciona como herramientas especialmente útiles cuando las necesidades incluyen esta capacidad para el preprocesamiento o modelado del problema. BHAVRP es la única biblioteca específicamente enfocada en resolver la fase de asignación del MDVRP, mediante la implementación de heurísticas clásicas y algoritmos de agrupamiento. Las demás bibliotecas analizadas abordan el problema completo o incluyen la asignación como un paso preliminar, sin especialización explícita en esta etapa.

BHAVRP actualmente ha sido integrado como parte funcional en herramientas como: CaSVRP y BHCVRP. Aunque su uso se ha concentrado en estos sistemas, su arquitectura modular y sus implementaciones en Java y Python permiten su potencial integración en otras aplicaciones compatibles con estos entornos \parencite{ericramosaragonIncorporacionModuloMetricas2025}.

Como se puede apreciar en la Figura~\ref{fig:flujo}, el flujo de ejecución comienza en el software que consume BHAVRP, donde se introducen los datos correspondientes a una instancia del problema MDVRP. Estos datos se envían al componente BHAVRP, el cual procede a cargar los datos del problema y a determinar el tipo de distancia que se utilizará en el cálculo de costos. A continuación, BHAVRP verifica si los datos fueron cargados correctamente. Si los datos presentan errores, el proceso se detiene. En caso contrario, se configuran los parámetros de ejecución necesarios, y se genera la matriz de costo correspondiente según la distancia seleccionada. Una vez generada la matriz de costos, el algoritmo se ejecuta en el contexto de BHAVRP, utilizando los parámetros definidos y la matriz de costos generada previamente. Al finalizar la ejecución del algoritmo, BHAVRP genera una solución, que es devuelta al sistema solicitante. Este último obtiene así los resultados de la asignación de clientes a depósitos, completando el ciclo de interacción  \parencite{ericramosaragonIncorporacionModuloMetricas2025}.

\figura{fig06_flujo_bhavrp.png}{1\textwidth}{Diagrama de actividades del flujo de BHAVRP.}{fig:flujo}

Para soportar la variedad de algoritmos y garantizar su correcta ejecución, BHAVRP presenta una arquitectura n-capas con un enfoque basado en reutilización, promoviendo consistencia con los lenguajes de desarrollo seleccionados  \parencite{ericramosaragonIncorporacionModuloMetricas2025}. Esta arquitectura se presenta en la Figura~\ref{fig:arquitectura}, que ofrece una vista detallada de BHAVRP.

\figura{fig07_arquitectura.png}{1\textwidth}{Diagrama de arquitectura n-capas con enfoque en reutilización de BHAVRP.}{fig:arquitectura}

La capa específica contiene los elementos responsables de gestionar la ejecución del sistema, como la clase Controller, que coordina la carga de datos, la configuración de parámetros y el flujo general de ejecución. Por su parte, en la capa general, que representa el núcleo de la biblioteca es responsable de la definición del problema, la modelación de las instancias y la ejecución de los algoritmos de asignación. Por su parte, la biblioteca trabaja con una estructura de datos que representan los elementos clave del MDVRP: clientes, depósitos, y sus respectivas restricciones y parámetros. La estructura básica para utilizar la biblioteca incluye un encabezado con el número de depósitos, clientes y vehículos disponibles por depósito, seguido de las capacidades máximas de los vehículos. Los clientes y los depósitos están definidos por un identificador y sus coordenadas (x, y). Además, cada cliente tiene una demanda asociada, mientras que cada depósito cuenta con una capacidad equivalente a la suma de las capacidades de los vehículos en su flota. BHAVRP en la capa intermedia contiene varias herramientas. Una de estas permite trabajar tanto con distancias reales como aproximadas (Euclidean, Manhattan, Chebyshev y Haversine). Luego para abordar el cálculo aproximado de las distancias, integra la herramienta OSRM (Open Source Routing Machine), que permite calcular distancias y tiempos de viaje basados en mapas reales, tanto en modalidad local como a través de solicitudes a un servidor en línea. En la capa de software se implementa la comunicación de la herramienta OSRM que expone sus servicios a través de una API HTTP RESTful. El cliente Java (como osrm-java-client) construye las peticiones HTTP y convierte las respuestas a JSON \parencite{ericramosaragonIncorporacionModuloMetricas2025}.

\subsection{El problema del ruido en los algoritmos de BHAVRP}
\label{sec:ruido_mdvrp}

La fase de asignación de clientes a depósitos en el MDVRP se ha abordado dentro de la biblioteca principalmente mediante heurísticas clásicas de asignación y algoritmos de agrupamiento particionales. Sin embargo, estos métodos presentan la limitación de que no están diseñados para detectar ni aislar puntos atípicos. Como resultado, generan clústeres que incluyen clientes alejados del depósito asignado lo que incrementa el costo de las rutas posteriores.

Esta limitación se evidencia en los resultados experimentales reportados en \parencite{ericramosaragonIncorporacionModuloMetricas2025} sobre una instancia real de 154 clientes y 5 depósitos en La Habana. La Figura~\ref{fig:ruido_tcc} muestra la asignación generada por la heurística \textit{Three Criteria Clustering} (TCC), donde se observa que algunos clientes asignados a un depósito se encuentran en zonas alejadas de su depósito principal. De forma similar, la Figura~\ref{fig:ruido_ff} presenta la asignación del algoritmo \textit{Farthest-First} (FF), caracterizada por clústeres con más dispersión y solapamiento entre zonas de distintos depósitos.

\figura{fig27_tcc.png}{0.85\textwidth}{Asignación generada por el algoritmo TCC sobre una instancia real.}{fig:ruido_tcc}

\figura{fig28_ff.png}{0.85\textwidth}{Asignación generada por el algoritmo FF sobre una instancia real.}{fig:ruido_ff}

Esto se ve reflejado en las métricas internas de validación. Algoritmos con más ruido obtienen los peores valores en métricas de cohesión y separación, lo que confirma que la presencia de este degrada la calidad estructural de la partición \parencite{ericramosaragonIncorporacionModuloMetricas2025}. Adicionalmente, el análisis de correlación entre las métricas internas y el costo total de las rutas evidencia que una asignación deficiente no siempre se traduce en un costo mayor.

Estos resultados fundamentan la necesidad de incorporar a BHAVRP un algoritmo de agrupamiento basado en densidad, capaz de identificar regiones de alta concentración de clientes y aislar aquellos puntos que no pertenecen a ninguna estructura densa. A diferencia de los métodos particionales y las heurísticas clásicas, los algoritmos basados en densidad como DBSCAN, HDBSCAN u OPTICS incorporan explícitamente un mecanismo de detección de ruido, lo que permite tratarlos de forma diferenciada en la fase de asignación. Aunque la biblioteca ya contiene una adaptación de DBSCAN este presenta ciertas limitaciones a considerar, por lo que es necesario implementar otro algoritmo de esta categoría.

\section{Conclusiones parciales}

A partir del estudio realizado sobre el marco teórico de los problemas de planificación de rutas de vehículos, específicamente la variante con múltiples depósitos, así como los métodos de solución y algoritmos de Minería de Datos existentes para resolver estos problemas se puede concluir que:

\begin{itemize}
  \item Los VRP son problemas de optimización combinatoria clasificados como NP-Duro por lo que se utilizan principalmente algoritmos heurísticos y metaheurísticos para su solución. El MDVRP es una de las numerosas variantes del VRP y requiere otro tipo de estrategias de solución.
  \item Una de las estrategias para la solución del MDVRP es dividir en la fase de asignación y la fase de planificación de las rutas. Además, para esta primera fase de asignación resulta efectivo el uso de algoritmos de minería de datos.
  \item BHAVRP es una biblioteca que permite resolver esa primera fase de asignación y tiene la facilidad para a\~nadir nuevas heurísticas y algoritmos de agrupamiento gracias a su arquitectura n-capas orientada a reutilización.
  \item La incorporación de otro algoritmo basado en densidad a BHAVRP enriquece la biblioteca con otro algoritmo de esta categoría que permita tener un mecanismo de detección de ruido.
\end{itemize}